% واژه‌های قیدی و پیوندی افزوده‌شده به گنجینهٔ واژه‌های فارسی
% برابرها همواره وابسته به بافت هستند؛ جدول‌ها جایگزینی خودکار نیستند.
\chapter{قیدها و واژه‌های پیوندی}
\noindent
برابر فارسی را باید بر پایهٔ معنی جمله برگزید؛ یک واژه در همهٔ بافت‌ها یک برابر ثابت ندارد.
برای نمونه، «دقیقاً پنج نفر» را می‌توان «درست پنج نفر» نوشت، اما «دقیقاً اجرا کن»
بهتر است «مو‌به‌مو اجرا کن» بازنویسی شود.
همچنین «شاید» از «به‌گمان زیاد» اطمینان کمتری را می‌رساند.

\section*{۱. درستی، گمان و تأکید}
\begingroup
\small
\setlength{\tabcolsep}{4pt}
\renewcommand{\arraystretch}{1.3}
\begin{longtable}{|p{2.5cm}|p{4.0cm}|p{2.3cm}|p{5.3cm}|}
\hline
واژهٔ رایج & برابر فارسی & انگلیسی & توضیح کاربرد \\
\hline
\endfirsthead
\hline
واژهٔ رایج & برابر فارسی & انگلیسی & توضیح کاربرد \\
\hline
\endhead

دقیقاً & درست / مو‌به‌مو / بی‌کم‌وکاست & \textenglish{Exactly} & برای شمار: «درست پنج نفر»؛ برای پیروی از دستور: «مو‌به‌مو». \\
\hline
حقیقتاً & به‌راستی / در حقیقت & \textenglish{Truly} & برای بیان راستی؛ در بسیاری از جمله‌ها «به‌راستی» روان‌تر است. \\
\hline
واقعاً & به‌راستی / به‌درستی & \textenglish{Really} & برای شگفتی یا تأکید؛ «به‌درستی» گاهی به معنی به‌شیوهٔ درست است. \\
\hline
یقیناً & بی‌گمان / بی‌تردید & \textenglish{Certainly} & برای اطمینان فراوان، نه پیش‌بینی نامطمئن. \\
\hline
مسلماً & بی‌گمان / بی‌تردید & \textenglish{Certainly} & برای تأکید بر امری دانسته یا پذیرفته‌شده. \\
\hline
احتمالاً & به‌گمان زیاد / شاید & \textenglish{Probably} & «شاید» از «به‌گمان زیاد» کم‌قطعیت‌تر است. \\
\hline
محتملاً & شاید / به‌گمان زیاد & \textenglish{Possibly} & بسته به میزان اطمینان نویسنده. \\
\hline
احیاناً & شاید / چه‌بسا & \textenglish{Possibly} & در پرسش‌های مؤدبانه نیز کاربرد دارد: «شاید بدانید؟». \\
\hline
ظاهراً & گویا / چنین پیداست & \textenglish{Apparently} & برداشت از نشانه‌هاست، نه اثبات قطعی. \\
\hline
قاعدتاً & بر پایهٔ روال / چنان‌که انتظار می‌رود & \textenglish{Normally} & برای نتیجهٔ مورد انتظار از یک قاعده یا روال. \\
\hline
مطلقاً & یکسره / به‌هیچ‌روی & \textenglish{Absolutely} & در نفی: «به‌هیچ‌روی»؛ در تأکید مثبت: «یکسره». \\
\hline
اساساً & از پایه / از بن & \textenglish{Fundamentally} & برای سخن گفتن از بنیاد موضوع. \\
\hline
اصولاً & در اصل / از بنیاد & \textenglish{In principle} & برای بیان قاعدهٔ کلی؛ همیشه با «اصلاً» یکسان نیست. \\
\hline
لزوماً & به‌ناگزیر / ناچار & \textenglish{Necessarily} & در ترکیب «نه لزوماً» بهتر است «نه همیشه» یا «نه به‌ناچار» با توجه به معنا بیاید. \\
\hline
\end{longtable}
\endgroup


\section*{۲. زمان، توالی و تکرار}
\begingroup
\small
\setlength{\tabcolsep}{4pt}
\renewcommand{\arraystretch}{1.3}
\begin{longtable}{|p{2.5cm}|p{4.0cm}|p{2.3cm}|p{5.3cm}|}
\hline
واژهٔ رایج & برابر فارسی & انگلیسی & توضیح کاربرد \\
\hline
\endfirsthead
\hline
واژهٔ رایج & برابر فارسی & انگلیسی & توضیح کاربرد \\
\hline
\endhead

جدیداً & به‌تازگی & \textenglish{Recently} & در اشاره به رخداد یا تغییر تازه. \\
\hline
اخیراً & به‌تازگی / همین چندی پیش & \textenglish{Recently} & برای رویدادی در گذشتهٔ نزدیک. \\
\hline
قبلاً & پیش‌تر / از پیش & \textenglish{Previously} & «از پیش» بیشتر برای آمادگی یا آگاهی پیشینی. \\
\hline
بعداً & سپس / پس از آن / دیرتر & \textenglish{Later} & «سپس» برای ترتیب؛ «دیرتر» برای زمان آینده. \\
\hline
فعلاً & اکنون / تا آن هنگام / تا اطلاع بعدی & \textenglish{For now} & زمان پایان یا موقتی‌بودن را از بافت جمله دریابید. \\
\hline
فوراً & بی‌درنگ / درجا & \textenglish{Immediately} & برای انجام کار بدون درنگ. \\
\hline
سریعاً & با شتاب / شتابان & \textenglish{Quickly} & سرعت را می‌رساند؛ نه لزوماً آغاز بی‌درنگ. \\
\hline
تدریجاً & رفته‌رفته / کم‌کم & \textenglish{Gradually} & دگرگونی در گذر زمان. \\
\hline
مجدداً & دوباره / از نو & \textenglish{Again} & «از نو» می‌تواند آغاز از ابتدا را برساند. \\
\hline
مکرراً & بارها / چندین‌بار & \textenglish{Repeatedly} & تکرار چندباره؛ نه لزوماً بی‌وقفه. \\
\hline
مرتباً & پی‌درپی / به‌طور منظم & \textenglish{Regularly} & «پی‌درپی» برای توالی؛ «منظم» برای نظم زمانی. \\
\hline
دائماً & پیوسته / همواره & \textenglish{Constantly} & برای استمرار، نه تکرار گاه‌به‌گاه. \\
\hline
موقتاً & برای چندی / کوتاه‌مدت & \textenglish{Temporarily} & با پایان احتمالی یا مشخص. \\
\hline
نهایتاً & سرانجام / در پایان & \textenglish{Ultimately} & برای فرجام یک روند؛ نه لزوماً بیشینهٔ مقدار. \\
\hline
ابتدائاً & نخست / در آغاز & \textenglish{Initially} & برای نخستین گام یا وضعیت آغازین. \\
\hline
متعاقباً & پس از آن / در پی آن & \textenglish{Subsequently} & ترتیب زمانی؛ گاه همراه با پیوند علّی. \\
\hline
متناوباً & یک‌درمیان / نوبتی & \textenglish{Alternately} & برای تناوب؛ با «پیوسته» تفاوت دارد. \\
\hline
سابقاً & پیش‌تر / در گذشته & \textenglish{Formerly} & برای وضعیت یا روال پیشین. \\
\hline
حالياً & اکنون / اینک & \textenglish{Currently} & برای زمان کنونی. \\
\hline
آنياً & همان دم / درجا & \textenglish{Instantly} & برای رخداد بی‌فاصله. \\
\hline
\end{longtable}
\endgroup


\section*{۳. شیوه، قصد و بیان}
\begingroup
\small
\setlength{\tabcolsep}{4pt}
\renewcommand{\arraystretch}{1.3}
\begin{longtable}{|p{2.5cm}|p{4.0cm}|p{2.3cm}|p{5.3cm}|}
\hline
واژهٔ رایج & برابر فارسی & انگلیسی & توضیح کاربرد \\
\hline
\endfirsthead
\hline
واژهٔ رایج & برابر فارسی & انگلیسی & توضیح کاربرد \\
\hline
\endhead

صریحاً & آشکارا / بی‌پرده & \textenglish{Explicitly} & برای بیان روشن و بی‌ابهام. \\
\hline
مستقیماً & یکراست / بی‌میانجی & \textenglish{Directly} & برای مسیر یا ارتباط بدون واسطه. \\
\hline
علناً & آشکارا / پیش چشم همگان & \textenglish{Publicly} & برای کاری که پنهانی نیست. \\
\hline
شخصاً & خودم / به‌شخصه & \textenglish{Personally} & در متن یکم‌شخص «خودم» روان‌تر است. \\
\hline
عمداً & آگاهانه / از روی خواست & \textenglish{Intentionally} & برای کار از روی قصد، نه اتفاق. \\
\hline
سهواً & ناخواسته / از روی اشتباه & \textenglish{Inadvertently} & نبود قصد را می‌رساند؛ لزوماً نبود آگاهی نیست. \\
\hline
تصادفاً & اتفاقی / از سر اتفاق & \textenglish{By chance} & برای رخداد پیش‌بینی‌نشده. \\
\hline
اتفاقاً & از قضا / دست‌برقضا & \textenglish{As it happens} & بسته به جمله می‌تواند هم‌زمانی یا غافلگیری را برساند. \\
\hline
اجباراً & به‌ناچار / از سر ناچاری & \textenglish{Under compulsion} & برای کاری بدون اختیار واقعی. \\
\hline
اختیاراً & به خواست خود / خودخواسته & \textenglish{Voluntarily} & برای کاری از روی گزینش آزادانه. \\
\hline
طبیعتاً & به‌خودی‌خود / چنان‌که پیداست & \textenglish{Naturally} & وابسته به این‌که منظور روند طبیعی باشد یا انتظار منطقی. \\
\hline
عیناً & مو‌به‌مو / بی‌کم‌وکاست & \textenglish{Verbatim} & برای بازگویی یا بازسازی بدون دگرگونی. \\
\hline
عملاً & در عمل / در کردار & \textenglish{In practice} & در برابر نظریه یا ادعا. \\
\hline
شفاهاً & زبانی / گفتاری & \textenglish{Orally} & گفته می‌شود، نوشته نمی‌شود. \\
\hline
کتباً & نوشتاری / در نوشته & \textenglish{In writing} & برای ثبت رسمی یا غیررسمی در نوشته. \\
\hline
اصالتاً & در اصل / از ریشه & \textenglish{Originally} & در زمینهٔ خاستگاه؛ با «ذاتاً» یکی نیست. \\
\hline
ذاتاً & در سرشت / سرشتاً & \textenglish{Inherently} & برای ویژگی درونی و بنیادین. \\
\hline
مستقلاً & جداگانه / به‌تنهایی & \textenglish{Independently} & بدون وابستگی؛ نه لزوماً در تنهایی فیزیکی. \\
\hline
تلویحاً & سربسته / به‌کنایه & \textenglish{Implicitly} & برای اشارهٔ غیرمستقیم. \\
\hline
توأماً & با هم / هم‌زمان & \textenglish{Together} & برای همراهی دو یا چند چیز. \\
\hline
متقابلاً & دوسویه / در برابر & \textenglish{Mutually} & برای کنش و واکنش دو سویه. \\
\hline
متأسفانه & شوربختانه / دریغا & \textenglish{Unfortunately} & برای افسوس یا ناخشنودی. \\
\hline
\end{longtable}
\endgroup


\section*{۴. اندازه، گستره و بسامد}
\begingroup
\small
\setlength{\tabcolsep}{4pt}
\renewcommand{\arraystretch}{1.3}
\begin{longtable}{|p{2.5cm}|p{4.0cm}|p{2.3cm}|p{5.3cm}|}
\hline
واژهٔ رایج & برابر فارسی & انگلیسی & توضیح کاربرد \\
\hline
\endfirsthead
\hline
واژهٔ رایج & برابر فارسی & انگلیسی & توضیح کاربرد \\
\hline
\endhead

کلاً & روی‌هم‌رفته / در مجموع & \textenglish{Overall} & جمع‌بندی؛ در «کلاً ممنوع است» می‌تواند «یکسره» باشد. \\
\hline
کاملاً & به‌تمامی / یکسره & \textenglish{Completely} & برای کامل‌بودن، نه نزدیک‌بودن. \\
\hline
تماماً & یکسره / به‌تمامی & \textenglish{Entirely} & برای همهٔ بخش‌ها. \\
\hline
اکثراً & بیشترشان / بیشتر وقت‌ها & \textenglish{Mostly} & بسته به اینکه شمار افراد باشد یا بسامد زمانی. \\
\hline
غالباً & بیشتر وقت‌ها / بیشتر & \textenglish{Often} & بسامد یا فراوانی را از بافت تشخیص دهید. \\
\hline
عموماً & به‌طور کلی / بیشتر وقت‌ها & \textenglish{Generally} & برای گزارهٔ عمومی، نه همیشه و همه‌جا. \\
\hline
معمولاً & بیشتر وقت‌ها / بنا بر روال & \textenglish{Usually} & برای چیزی که روال رایج است. \\
\hline
نسبتاً & تا اندازه‌ای / کمابیش & \textenglish{Relatively} & در قیاس با چیزی دیگر. \\
\hline
شدیداً & سخت / بسیار & \textenglish{Intensely} & برای شدت؛ «سخت» همواره مناسب متن علمی نیست. \\
\hline
عمیقاً & ژرف / از ژرفای دل & \textenglish{Deeply} & برای ژرفای معنا یا احساس. \\
\hline
صرفاً & تنها / تنها و بس & \textenglish{Merely} & برای انحصار یا محدودکردن دامنه. \\
\hline
منحصراً & تنها / ویژهٔ & \textenglish{Exclusively} & برای انحصار؛ «ویژهٔ» در ترکیب اسمی. \\
\hline
خصوصاً & به‌ویژه & \textenglish{Especially} & برای برجسته‌کردن یک مورد. \\
\hline
مخصوصاً & به‌ویژه / ویژه & \textenglish{Especially} & بسته به ساختار جمله. \\
\hline
بالاخص & به‌ویژه / بیش از همه & \textenglish{In particular} & برای تأکید بر مصداق مشخص. \\
\hline
اجمالاً & به‌کوتاهی / کوتاه سخن & \textenglish{Briefly} & برای گزارش کوتاه و فشرده. \\
\hline
تفصیلاً & با ریزه‌کاری / به‌گونهٔ گسترده & \textenglish{In detail} & برای شرح کامل؛ نه تنها بیان طولانی. \\
\hline
\end{longtable}
\endgroup


\section*{۵. پیوند و سازمان‌دهی گفتار}
\begingroup
\small
\setlength{\tabcolsep}{4pt}
\renewcommand{\arraystretch}{1.3}
\begin{longtable}{|p{2.5cm}|p{4.0cm}|p{2.3cm}|p{5.3cm}|}
\hline
واژهٔ رایج & برابر فارسی & انگلیسی & توضیح کاربرد \\
\hline
\endfirsthead
\hline
واژهٔ رایج & برابر فارسی & انگلیسی & توضیح کاربرد \\
\hline
\endhead

ضمناً & ناگفته نماند / افزون بر این & \textenglish{Incidentally} & برای افزودن نکته؛ «همچنین» در برخی جمله‌ها بهتر است. \\
\hline
مضافاً & افزون بر این / گذشته از این & \textenglish{Additionally} & برای آوردن دلیل یا نکتهٔ دیگر. \\
\hline
بعلاوه & افزون بر این / گذشته از آن & \textenglish{Moreover} & برای افزودن نکته. \\
\hline
اولاً & نخست / نخست اینکه & \textenglish{Firstly} & برای آغاز یک ترتیب. \\
\hline
ثانیاً & دوم / دوم اینکه & \textenglish{Secondly} & برای گام دوم در استدلال. \\
\hline
ثالثاً & سوم / سوم اینکه & \textenglish{Thirdly} & برای گام سوم در استدلال. \\
\hline
اختصاراً & به‌کوتاهی / کوتاه بگویم & \textenglish{In brief} & برای خلاصه‌گویی. \\
\hline
نتیجتاً & از این رو / در نتیجه & \textenglish{Consequently} & برای پیامد یک دلیل؛ «از این رو» روان‌تر است. \\
\hline
بالعکس & برعکس / وارونه & \textenglish{Conversely} & بسته به اینکه تقابل باشد یا وارونگی. \\
\hline
فی‌الواقع & به‌راستی / در حقیقت & \textenglish{In fact} & برای روشن‌کردن واقعیت. \\
\hline
فی‌نفسه & به‌خودی‌خود & \textenglish{In itself} & بدون درنظرگرفتن اثرهای بیرونی. \\
\hline
علی‌رغم & با همهٔ / با اینکه & \textenglish{Despite} & «با اینکه» برای پیوند دو گزاره. \\
\hline
بالنتیجه & از این رو / در پی آن & \textenglish{Therefore} & برای نشان‌دادن نتیجه. \\
\hline
مع‌هذا & بااین‌همه / بااین‌حال & \textenglish{Nevertheless} & برای بیان تقابل با جملهٔ پیشین. \\
\hline
\end{longtable}
\endgroup


\section*{چند نمونهٔ بازنویسی}
\begin{itemize}
\item «احتمالاً این مشکل مجدداً رخ می‌دهد» $\leftarrow$ «به‌گمان زیاد این مشکل دوباره پیش می‌آید».
\item «قطعاً باید مستقیماً گفت‌وگو کنیم» $\leftarrow$ «بی‌گمان باید یکراست گفت‌وگو کنیم».
\item «اخیراً این خطا مکرراً دیده شده است» $\leftarrow$ «این خطا به‌تازگی بارها دیده شده است».
\end{itemize}

\noindent
\textbf{یادآوری:} واژه‌هایی مانند «قطعاً»، «حتماً»، «اصلاً»، «تقریباً»، «مثلاً» و «عمدتاً»
در بخش‌های پیشین آمده‌اند و برای پرهیز از مدخل‌های تکراری اینجا دوباره فهرست نشده‌اند.
